\chapter{System Architecture}

\section{The stage pipeline}

A run is one YAML file (\texttt{configs/*.yaml}) interpreted by a chain of
pure, testable stages. Each stage reads only the sub-model of \texttt{RunConfig}
it needs, so configs stay flat and reviewable, and each stage can be tested in
isolation with synthetic inputs.

\begin{figure}[h]
\centering
\begin{tikzpicture}[node distance=7mm]
  \node[stage] (cfg)  {config};
  \node[stage, right=of cfg]  (dat)  {data};
  \node[stage, right=of dat]  (feat) {features};
  \node[stage, right=of feat] (sig)  {signals};
  \node[stage, right=of sig]  (val)  {validation};
  \draw[flowarr] (cfg)--(dat); \draw[flowarr] (dat)--(feat);
  \draw[flowarr] (feat)--(sig); \draw[flowarr] (sig)--(val);

  \node[stage, below=13mm of val] (bt)   {backtest};
  \node[stage, left=of bt]  (risk) {risk};
  \node[stage, left=of risk] (exec) {execution};
  \node[stage, left=of exec] (rep)  {report};
  \node[stage, left=of rep]  (live) {live loop};
  \draw[flowarr] (val.south) -- (bt.north);
  \draw[flowarr] (bt)--(risk); \draw[flowarr] (risk)--(exec);
  \draw[flowarr] (exec)--(rep); \draw[flowarr] (rep)--(live);

  \node[stagealt, below=13mm of live] (mcp) {mcp/server.py};
  \draw[flowarr] (live.south) -- (mcp.north);
  \node[note, right=4mm of mcp, text width=6.4cm] {exposes \texttt{run\_backtest},
    \texttt{paper\_iteration}, \texttt{list\_runs}, \ldots{} to any MCP client};
\end{tikzpicture}
\caption{The eleven modules of \texttt{src/qp/}, in call order. The bottom row
reads right to left to keep the diagram compact; data still flows
backtest~$\to$~risk~$\to$~execution~$\to$~report~$\to$~live.}
\end{figure}

\section{Orchestration}

\texttt{qp.pipeline.run\_pipeline(cfg)} is the single entry point used by the
CLI's \texttt{demo} and \texttt{backtest} commands:

\begin{lstlisting}
long_df = load_prices(cfg.data)                  # data
close   = to_wide(long_df, "close")
signal  = make_signal(cfg.signal, cfg.features, cfg.validation)
scores  = signal.generate(close)                 # features + signals (fused)
result  = run_backtest(close, scores, cfg)        # validation + backtest + risk
write_tearsheet(result, cfg.run_dir())           # report
\end{lstlisting}

\texttt{qp.live.loop.run\_once(cfg)} reuses the same \texttt{data}, \texttt{features},
and \texttt{signals} code, then diverges into \texttt{execution} instead of
\texttt{backtest} -- the live loop and the backtester are guaranteed to compute
the same target weights from the same signal, because they call the same
functions.

\section{Why this shape}

\begin{designnote}
Every stage is a pure function or a small class with one public method
(\texttt{generate}, \texttt{fetch}, \texttt{rebalance\_to}). Nothing reaches into
another stage's internals. That buys three things directly relevant to career
capital: (1) each stage is unit-testable against synthetic data without a
network or a broker; (2) a new signal, data source, or broker is additive, not
a rewrite; (3) the same code path serves the backtester and the live loop, so
there is no ``it worked in backtest but not live'' class of bug from divergent
implementations.
\end{designnote}
